\section{CONSIDERAÇÕES FINAIS}

O desenvolvimento deste Trabalho de Conclusão de Curso resultou na concepção, implementação e validação da plataforma WeFIND, uma aplicação móvel colaborativa planejada para conectar tutores e voluntários na localização de animais perdidos e no incentivo à posse responsável. O ponto de partida da pesquisa comprovou que as alternativas habitualmente acionadas pela população, como cartazes impressos afixados em postes e postagens dispersas em redes sociais generalistas, mostram-se ineficazes perante a urgência das buscas, pois carecem de mapas georreferenciados, perdem alcance rapidamente no fluxo de mensagens e expõem números particulares a riscos de fraude.

A avaliação do ciclo de desenvolvimento confirma o pleno atendimento dos objetivos traçados para o projeto. O diagnóstico inicial com tutores e protetores fundamentou a identificação das principais dores da comunidade, servindo de base empírica para o levantamento de 43 requisitos funcionais e 14 requisitos não funcionais. A análise comparativa de plataformas similares permitiu mapear limitações existentes no mercado e direcionar o WeFIND como uma solução gratuita, dotada de mapa interativo por GPS e canal de diálogo privativo. 

Em termos técnicos, a modelagem conceitual formalizou as regras de negócio em diagramas da UML, e a infraestrutura foi viabilizada por meio da biblioteca React Native, do ecossistema Expo e dos serviços em nuvem do Supabase com banco de dados relacional PostgreSQL. O design centrado no usuário assegurou clareza tipográfica com a fonte Roboto e coerência na paleta cromática, culminando em uma interface amigável. A validação prática com 20 voluntários ratificou o sucesso da proposta, alcançando a pontuação média de 85,3 pontos na escala System Usability Scale, situando o WeFIND no patamar de usabilidade excelente.

\subsection{Dificuldades Encontradas}

Durante o desenvolvimento do projeto, foram enfrentados desafios técnicos relevantes. A renderização simultânea de múltiplos marcadores em dispositivos com capacidade de processamento mais modesta exigiu a otimização de consultas espaciais e a reorganização de camadas do mapa para evitar travamentos visuais. No tratamento de arquivos multimídia, foi necessário implementar rotinas de compressão e redimensionamento de imagens diretamente no aparelho do usuário antes da transmissão ao Supabase Storage, economizando dados de internet móvel. 

No que tange à comunicação em tempo real, demandou-se esforço para sincronizar a entrega de mensagens instantâneas garantindo o isolamento estrito de dados por meio de políticas de segurança em nível de linha (RLS) no banco de dados. Ademais, realizou-se uma exploração preliminar com modelos de inteligência artificial para moderação automatizada de fotos, buscando garantir que apenas fotos de animais fossem submetidas. Contudo, os modelos testados apresentaram oscilações de precisão em condições de iluminação variável ou enquadramentos parciais, o que poderia bloquear envios legítimos em momentos críticos de emergência. Em virtude disso, optou-se por priorizar a moderação assistida pela comunidade por meio de denúncias e análise pelo administrador, preservando a agilidade dos tutores.

\subsection{Trabalhos Futuros}

Como continuidade à evolução da plataforma WeFIND, sugerem-se oportunidades para desdobramentos futuros. Recomenda-se a pesquisa e integração de modelos aprimorados de visão computacional para triagem automática de imagens de animais, bem como o estudo de algoritmos de reconhecimento morfológico capazes de cruzar automaticamente as características visuais de animais perdidos com anúncios de animais recém-encontrados.

Propõe-se também a viabilização física do Perfil Digital do Pet, permitindo a gravação do código QR exclusivo em plaquetas metálicas de identificação acopláveis a coleiras tradicionais, o que facilitaria a identificação imediata por qualquer pedestre mesmo sem conexão ativa no momento do resgate. Por fim, sugere-se o desenvolvimento de um módulo específico para clínicas veterinárias e estabelecimentos parceiros, estabelecendo murais digitais comunitários de apoio à causa animal e expandindo a rede solidária de reencontros.

Dessa forma, o WeFIND cumpre seu propósito ao demonstrar que a integração entre engenharia de software móvel e solidariedade comunitária constitui um caminho viável e transformador, oferecendo uma ferramenta tecnológica acessível para abreviar a angústia da perda e promover o reencontro entre tutores e seus animais de estimação.
\clearpage
